\chapter{Appendix}

\section{API Documentation}
\subsection{Complete API Endpoints}

\subsubsection{Health Check}
\begin{lstlisting}[language=bash, caption=Health Check Endpoint]
GET /health

Response:
{
    "status": "healthy",
    "version": "2.0.0",
    "model_loaded": true,
    "symptoms_count": 132,
    "diseases_count": 41,
    "features": [
        "Disease Prediction",
        "NLP Symptom Extraction",
        "Severity Detection",
        "First Aid Guidance",
        "Conversational AI Agent"
    ]
}
\end{lstlisting}

\subsubsection{Predict from Symptoms}
\begin{lstlisting}[language=bash, caption=Predict from Symptoms]
POST /predict/symptoms
Content-Type: application/json

Request Body:
{
    "symptoms": ["fever", "headache", "fatigue"],
    "top_k": 5,
    "duration": "3 days"
}

Response:
{
    "predictions": [
        {
            "disease": "Common Cold",
            "probability": 0.85,
            "confidence_percent": "85.0%",
            "description": "Viral infection...",
            "precautions": ["rest", "hydration"],
            "severity_score": 2.5,
            "severity_level": "MILD"
        }
    ],
    "matched_symptoms": ["high_fever", "headache", "fatigue"],
    "severity": {
        "level": "MILD",
        "score": 2.5,
        "confidence": 0.8,
        "reason": "Mild symptoms...",
        "is_emergency": false
    },
    "recommended_specialist": {
        "key": "general_practitioner",
        "name": "General Practitioner",
        "description": "...",
        "icon": "👨‍⚕️"
    }
}
\end{lstlisting}

\subsubsection{Predict from Text}
\begin{lstlisting}[language=bash, caption=Predict from Natural Language Text]
POST /predict/text
Content-Type: application/json

Request Body:
{
    "text": "I have had fever and headache for 2 days",
    "top_k": 5
}

Response: (Same structure as /predict/symptoms)
\end{lstlisting}

\subsubsection{Chat Endpoint}
\begin{lstlisting}[language=bash, caption=Conversational Agent]
POST /chat
Content-Type: application/json

Request Body:
{
    "session_id": "user123",
    "message": "I have fever and body pain"
}

Response:
{
    "response": "I've noted: Fever, Body Pain...",
    "state": "collecting_symptoms",
    "suggestions": ["I also have chills", "That's all"],
    "collected_symptoms": ["high_fever", "muscle_pain"],
    "predictions": null,
    "severity": null,
    "is_emergency": false
}
\end{lstlisting}

\section{System Configuration}
\subsection{Configuration File}
\begin{lstlisting}[language=Python, caption=System Configuration]
# Model Configuration
MODEL_DIR = Path("models")
RANDOM_STATE = 42
TEST_SIZE = 0.2
TOP_K_DISEASES = 5

# Dataset Paths
DATASET_PATH = "dataset.csv"
SYMPTOM_DESCRIPTION_PATH = "symptom_Description.csv"
SYMPTOM_PRECAUTION_PATH = "symptom_precaution.csv"
SYMPTOM_SEVERITY_PATH = "Symptom-severity.csv"

# API Configuration
API_HOST = "0.0.0.0"
API_PORT = 8000

# Severity Thresholds
SEVERITY_THRESHOLDS = {
    "MILD": 0.0,
    "MODERATE": 3.0,
    "SEVERE": 5.0,
    "EMERGENCY": 7.0
}
\end{lstlisting}

\section{Model Training Script}
\subsection{Training Code Example}
\begin{lstlisting}[language=Python, caption=Model Training Script]
from disease_predictor import DiseasePredictionModel

# Initialize model
model = DiseasePredictionModel()

# Load and preprocess data
df = model.preprocessor.load_main_dataset()
model.preprocessor.load_auxiliary_data()
X, y = model.preprocessor.prepare_training_data(df)

# Train model
model.train(X, y, use_ensemble=True)

# Save model
model.save()

# Test prediction
test_symptoms = ['fever', 'headache', 'fatigue']
predictions = model.predict(test_symptoms, top_k=3)
print(predictions)
\end{lstlisting}

\section{Dataset Statistics}
\subsection{Dataset Overview}
\begin{table}[H]
\centering
\caption{Dataset Statistics}
\begin{tabular}{|l|c|}
\hline
\textbf{Statistic} & \textbf{Value} \\
\hline
Total Records & 4,922 \\
\hline
Number of Diseases & 41 \\
\hline
Number of Symptoms & 132 \\
\hline
Average Symptoms per Disease & 5-8 \\
\hline
Most Common Disease & Common Cold \\
\hline
Dataset Format & CSV \\
\hline
\end{tabular}
\end{table}

\section{Supported Diseases}
The system supports prediction for the following 41 diseases:

\begin{itemize}
    \item Common Cold, Typhoid, Malaria, Dengue
    \item Diabetes, Hypertension, Heart Attack
    \item Gastroenteritis, Peptic Ulcer, GERD
    \item Migraine, Tension Headache, Cervical Spondylosis
    \item Arthritis, Osteoarthritis
    \item Fungal Infection, Allergy, Acne
    \item Urinary Tract Infection, Chronic Kidney Disease
    \item Jaundice, Hepatitis
    \item Pneumonia, Bronchial Asthma
    \item And 20+ other conditions
\end{itemize}

\section{Supported Symptoms}
The system recognizes 132 symptoms including:

\begin{itemize}
    \item Fever-related: high\_fever, mild\_fever, chills, sweating
    \item Pain-related: headache, stomach\_pain, chest\_pain, joint\_pain, muscle\_pain
    \item Respiratory: cough, breathlessness, phlegm, runny\_nose
    \item Digestive: nausea, vomiting, diarrhoea, constipation, acidity
    \item Skin: itching, skin\_rash, yellowish\_skin
    \item Neurological: dizziness, blurred\_vision, anxiety
    \item And 100+ other symptoms
\end{itemize}

\section{Emergency Contact Numbers}
The system includes the following emergency contact numbers for Pakistan:

\begin{table}[H]
\centering
\caption{Emergency Contact Numbers}
\begin{tabular}{|l|c|}
\hline
\textbf{Service} & \textbf{Number} \\
\hline
Rescue (Ambulance) & 1122 \\
\hline
Edhi Foundation & 115 \\
\hline
Police & 15 \\
\hline
Fire Department & 16 \\
\hline
Aman Foundation & 1021 \\
\hline
\end{tabular}
\end{table}

\section{Installation Instructions}
\subsection{Prerequisites}
\begin{enumerate}
    \item Python 3.8 or higher
    \item pip package manager
    \item 2GB+ RAM
    \item 1GB+ disk space
\end{enumerate}

\subsection{Installation Steps}
\begin{lstlisting}[language=bash, caption=Installation Commands]
# Clone repository
git clone <repository-url>
cd fyp/ai_agent

# Install dependencies
pip install -r requirements.txt

# Train models (first time only)
python train_improved.py

# Start API server
python api_v2.py
\end{lstlisting}

\section{Project Structure}
\begin{lstlisting}[language=bash, caption=Project Directory Structure]
ai_agent/
├── api_v2.py              # FastAPI REST endpoints
├── config.py              # Configuration settings
├── data_preprocessing.py  # Data loading and preprocessing
├── disease_predictor.py   # ML model training and prediction
├── symptom_checker_v2.py  # Conversational AI agent
├── nlp_processor.py       # NLP for symptom extraction
├── severity_model.py      # Severity detection model
├── first_aid.py           # First aid guidance system
├── specialist_mapper.py   # Specialist recommendations
├── maps_integration.py    # Location-based doctor search
├── prediction_validator.py # Prediction validation
├── duration_analyzer.py    # Duration-based analysis
├── models/                # Saved ML models
│   ├── ensemble_model.joblib
│   ├── rf_model.joblib
│   ├── xgb_model.joblib
│   ├── preprocessor.joblib
│   └── severity_model.joblib
├── web/                   # Web interface
│   └── index.html
└── requirements.txt       # Python dependencies
\end{lstlisting}

\section{Testing}
\subsection{Unit Tests}
The system includes comprehensive unit tests for:
\begin{itemize}
    \item Symptom extraction from text
    \item Disease prediction accuracy
    \item Severity detection
    \item First aid guide retrieval
    \item API endpoint functionality
\end{itemize}

\subsection{Running Tests}
\begin{lstlisting}[language=bash, caption=Test Execution]
# Run comprehensive tests
python test_comprehensive.py

# Test specific components
python -m pytest tests/
\end{lstlisting}

\section{Performance Benchmarks}
\subsection{Model Inference Times}
\begin{table}[H]
\centering
\caption{Inference Performance}
\begin{tabular}{|l|c|}
\hline
\textbf{Operation} & \textbf{Time (ms)} \\
\hline
Single Prediction & 50-100 \\
\hline
Top-5 Predictions & 80-150 \\
\hline
NLP Symptom Extraction & 20-50 \\
\hline
Severity Assessment & 30-60 \\
\hline
First Aid Lookup & 10-30 \\
\hline
\end{tabular}
\end{table}

\section{Error Codes}
The API uses standard HTTP status codes:
\begin{itemize}
    \item 200: Success
    \item 400: Bad Request (invalid input)
    \item 404: Not Found (resource not available)
    \item 500: Internal Server Error
    \item 503: Service Unavailable (model not loaded)
\end{itemize}

\section{License and Disclaimer}
\subsection{License}
This project is developed for academic purposes as part of Final Year Project at FAST-NUCES.

\subsection{Medical Disclaimer}
\textbf{IMPORTANT}: This system is for informational purposes only and does not provide medical advice, diagnosis, or treatment. Always consult qualified healthcare professionals for medical concerns. The system should not be used as a substitute for professional medical care, especially in emergency situations.

\section{Contact Information}
For questions or issues regarding this project:
\begin{itemize}
    \item Project Repository: [GitHub URL]
    \item Email: [Contact Email]
    \item Institution: FAST-NUCES, Islamabad
\end{itemize}
